\scanpage{84}
\diagram{\begin{tikzcd}[ampersand replacement=\&]
H \handwritten{\arrow[r,"\handwritten{u_\nu}"]} \& G_{\handwritten{\nu}}\ \handwritten{(} \handwritten{\arrow[r]} \& G\handwritten{)}\ ,
\end{tikzcd}}
lui-même déterminé par la connaissance de l'homomorphisme en sens inverse sur les groupes de caractères
\diagram{\begin{tikzcd}[ampersand replacement=\&]
Q \& \handwritten{\overline}{\handwritten{\underline}{Q}}^{\handwritten{*}} \handwritten{\arrow[l,"\handwritten{u_\nu^*}"']} \& \handwritten{\overline}{\handwritten{\underline}{Q}}^{\handwritten{*}} \handwritten{\arrow[l,"\handwritten{\lambda^\nu}"']} \handwritten{\arrow[ll,bend left=25,"\handwritten{|\lambda|}"']}\quad .
\end{tikzcd}}
On lit sur ce diagramme que, pour l'isomorphisme naturel du groupe des caractères de $G_{\handwritten{\nu}}$ avec $\handwritten{\overline}{\handwritten{\underline}{Q}}^{\handwritten{*}}$, provenant de l'élément $\handwritten{\underline}{f}_q\handwritten{\in}G_{\handwritten{\nu}}(\handwritten{\underline}{Q})$, \handwritten{on} \deleted{qui donne}\inserted{\handwritten{a}}
\[
u_{\handwritten{\nu}}^{\handwritten{*}}(\handwritten{\lambda})=v_p(\handwritten{\lambda})/\handwritten{\nu}\quad .
\]
On trouve donc\deleted{x} le résultat suivant de Manin, énoncé sous une forme plus jolie et débaréssé d'hypothèses restrictives parasites:

\handwritten{\uline}{Proposition} 10.3. Soit $(E,F_E)$ un F-isocristal sur $\unsure{\handwritten{\underline}{F}|\handwritten{\underline}{\mathbb{F}}}_q$, $q=p^{\handwritten{\nu}}$, et considérons le polynôme ca\unsure{\handwritten{r}|r}actéristique $f(T)$ de $F_E{}^{\handwritten{\nu}}$, \margin{\handwritten{NB si on prend plutôt la val. $\unsure{v_q|\nu_q}$ telle que $v_q(q)=1$, on trouve les valeurs propres $v_q(\lambda_i)$ (sans diviser par $\nu$)}}
soient $\handwritten{\lambda}_1,..,\handwritten{\lambda}_N$ ($N=\operatorname{rang}E$) ses racines, et $c_1$ , $..,c_{\handwritten{N}}$ leurs valeurs absolues normalisées ($c_i=v_p(\handwritten{\lambda}_i)$, $v_p(p)=1$). Alors \deleted{sur} dans $\operatorname{Rep}(\unsure{\handwritten{\underline}{H}|\handwritten{\underline}{G}})\handwritten{\otimes}_{\handwritten{\underline}{Q}_p}\handwritten{\overline}{\handwritten{\underline}{Q}}_p\handwritten{\simeq}\unsure{\operatorname{Fcriso}|\operatorname{Ecriso}}(\handwritten{\overline}{\handwritten{\underline}{F}}_q)\handwritten{\otimes}_{\handwritten{\underline}{Q}_p}\handwritten{\overline}{\handwritten{\underline}{Q}}_p$, l'image de $(E,E_E)$ a comme pentes les nombres $c_1/\handwritten{\nu}$, $...$ $c_N/\handwritten{\nu}$, de sorte que, si les \inserted{différentes} pentes possibles sont $\deleted{p_i}$\deleted{ avec multiplicité} $p_i=s_i/r_i$ ($s_i,r_i\handwritten{\in}\handwritten{\underline}{Z}\deleted{\unsure{x|r}}$, $r_i\handwritten{\geq}1$, $(s_i,r_i)=1$) \deleted{\handwritten{$\unsure{d_ip_i\in Q|d_i,p_i\in Q|d_ip_i\in Z|d_i,p_i\in Z}$}} avec multiplicité $d_i$, on a $r_i\handwritten{\mid}d_i$ i.e\handwritten{,} $d_i=d'_ir_i$, et l'image de $(E,F_E)$ dans $\operatorname{Fcriso}(\handwritten{\overline}{\handwritten{\underline}{F}}_p)$ est isomorphe à $\handwritten{\sum_i}E_{p_i}^{d'_i}$.


\handwritten{\uline{Cor.}}
% archive sheet number: 83

